\section{Formal Semantics}

This appendix provides precise definitions for the core semantic objects used throughout the manuscript. All definitions are structural rather than representational.

\subsection{System}

\begin{definition}[System]
A system is a tuple
\[
S \defeq (\States, \Constraints, \ActionFn)
\]
where:
\begin{itemize}
  \item $\States$ is an internal state space,
  \item $\Constraints : \States \to \{0,1\}$ is a constraint predicate,
  \item $\ActionFn : \States \to \Actions \cup \{\Refusal\}$ is an action function.
\end{itemize}
\end{definition}

\subsection{Agent}

\begin{definition}[Agent]
A system $S$ is an agent iff:
\[
\forall x \in \States,\quad \ConstraintFn(x) = 0 \Rightarrow \ActionFn(x) = \Refusal.
\]
That is, execution is conditionally suspended under constraint violation.
\end{definition}

\subsection{Didactic Interface}

\begin{definition}[Didactic Interface]
A system $D$ is a didactic interface relative to an execution environment $\Executable$ iff:
\[
\exists y \in \Actions \text{ such that } \Executable(y) \text{ succeeds independently of } \Constraints.
\]
\end{definition}

\subsection{Semantic Compression}

\begin{definition}[Semantic Compression]
A semantic compression is a map
\[
\AbstractFn : (\States, \Constraints) \to \Actions
\]
such that:
\[
\MutualInfo(\Constraints; \AbstractFn(\States)) < \MutualInfo(\Constraints; \States).
\]
\end{definition}

Constraint information loss is therefore structural, not accidental.

\subsection{Refusal}

\begin{definition}[Refusal]
Refusal is the distinguished output $\Refusal$ indicating suspension without substitution.
\end{definition}

Refusal is not an alternative action; it is the absence of execution.

